\section{Plan d'action dat\'e}\label{sec:plan}
% Chapitre 10 (tache 10, passage chapitres 09-10). Figure: fig_zeitplan (jalons du
% scenario de base, colonne "libelle" ecrite par kapitel_10). Macros: celles des chapitres
% 4, 5 et 8, plus _macros_redaction_kapitel_9_10 (jalons, annees des choix, epargnes,
% reserve d'urgence, revue avant l'objectif). Seul chapitre ou des listes a puces sont
% admises (plan d'action). Phrases liees aux donnees actuelles (a relire si les donnees
% changent): "un seul levier est puissant", "tout en maison ne fait gagner qu'une annee",
% "le capital de depart couvre une partie de la reserve".

Ce chapitre r\'epond \`a la question~: que dois-tu faire, et quand, \`a partir de septembre 2027~? Le
chapitre~\ref{sec:presse} s'est conclu sur un constat que la presse ne chiffre pas~: vivre des seuls
dividendes est confortable, mais co\^ute du capital et des ann\'ees de travail. La
section~\ref{sec:plan-verdict} dit franchement ce que le sc\'enario de base permet et quels choix
avanceraient ton d\'epart. La section~\ref{sec:plan-premiers} dresse la liste des premiers pas de
l'automne 2027, la section~\ref{sec:plan-jalons} fixe les jalons de ton capital, et la
section~\ref{sec:plan-revision} le calendrier des r\'evisions jusqu'\`a l'objectif.

\subsection{Ce que le sc\'enario de base permet, et ce qui l'avancerait}\label{sec:plan-verdict}

\bk{Au sc\'enario de base, r\'epartition \ReferenzStrategie{} et \ReferenzSparquote~\% d'\'epargne, tu
atteins l'objectif en \AnneeYannBasis, \`a \AlterYannBasis~ans~: c'est l'\^age l\'egal de \AgeLegal~ans.
Ce plan ne te donne donc, \`a ce taux, aucune retraite anticip\'ee.} \`A \SparquoteZwanzig~\%
d'\'epargne, aucune strat\'egie de dividendes n'atteint l'objectif avant \AlterHorizon~ans
(chapitre~\ref{sec:accumulation}). Le tableau~\ref{tab:choix-avancer-depart} rassemble les choix qui
changent cette date, avec l'\'epargne qu'ils demandent la premi\`ere ann\'ee.

\begin{table}[htbp]
  \centering\small
  \caption{\^Age auquel l'objectif de \ZielMonat~€ nets par mois est atteint selon la strat\'egie et le
  taux d'\'epargne, et \'epargne mensuelle de la premi\`ere ann\'ee (capital de d\'epart de
  \KapitalDepartMitte~€, euros de 2026).}\label{tab:choix-avancer-depart}
  \begin{tabularx}{\linewidth}{>{\bfseries\raggedright\arraybackslash}X
    *{3}{>{\centering\arraybackslash}p{0.17\linewidth}}}
    \toprule
    \rowcolor[gray]{0.9}
    \textbf{Strat\'egie} & \textbf{Taux d'\'epargne} & \textbf{\'Epargne par mois (€)}
      & \textbf{\^Age d'atteinte} \\
    \midrule
    Dividendes, \ReferenzStrategie & \SparquoteZwanzig~\% & \EpargneMoisZwanzig
      & \AlterMaisonSiebzigZwanzig \\
    Dividendes, \ReferenzStrategie{} (base) & \ReferenzSparquote~\% & \EpargneMoisBasis
      & \AlterYannBasis \\
    Dividendes, \ReferenzStrategie & \QuoteAnticipeeMaisonSiebzig~\%
      & \EpargneMoisAnticipeeMaisonSiebzig & \AlterAnticipeeMaisonSiebzig \\
    Dividendes, \ReferenzStrategie & \SparquoteFuenfzig~\% & \EpargneMoisFuenfzig
      & \AlterMaisonSiebzigFuenfzig \\
    Dividendes, tout en maison & \SparquoteDreissig~\% & \EpargneMoisDreissig
      & \AlterMaisonHundertDreissig \\
    ETF monde, retrait \TauxRetrait~\% & \QuoteAnticipeeEtfReferenz~\%
      & \EpargneMoisAnticipeeEtfReferenz & \AlterAnticipeeEtfReferenz \\
    ETF monde, retrait \TauxRetrait~\% & \SparquoteZwanzig~\% & \EpargneMoisZwanzig
      & \AlterEtfReferenzZwanzig \\
    ETF monde, retrait \TauxRetrait~\% & \SparquoteDreissig~\% & \EpargneMoisDreissig
      & \AlterEtfBasis \\
    ETF monde, retrait \TauxRetrait~\% & \SparquoteFuenfzig~\% & \EpargneMoisFuenfzig
      & \AlterEtfReferenzFuenfzig \\
    \bottomrule
  \end{tabularx}
\end{table}

Ce tableau montre deux leviers, et un seul est puissant. Le premier est le taux d'\'epargne~: en
restant fid\`ele aux dividendes, il faut \'epargner au moins \QuoteAnticipeeMaisonSiebzig~\% de ton
salaire net, \EpargneMoisAnticipeeMaisonSiebzig~€ par mois au d\'ebut, pour partir en
\AnneeAnticipeeMaisonSiebzig, \`a \AlterAnticipeeMaisonSiebzig~ans, et \SparquoteFuenfzig~\% pour
partir en \BrueckeAnneeDebutFuenfzig, \`a \AlterMaisonSiebzigFuenfzig~ans. Ce d\'epart anticip\'e a un
prix que le chapitre~\ref{sec:retraite} a chiffr\'e~: un pont de \BrueckeJahreFuenfzig~ans pendant
lequel tes dividendes sont ton seul revenu, et une rente l\'egale r\'eduite de
\RenteBrueckePerteProzent~\% \`a vie. Placer tout l'ETF maison en actions en direct ne fait gagner
qu'une ann\'ee au taux de base.

Le second levier est la strat\'egie. Un ETF monde capitalisant dont tu retires \TauxRetrait~\% par an
demande \KapitalEtfReferenz~€ au lieu de \KapitalMaisonSiebzig~€ (chapitre~\ref{sec:capital}) et
atteint l'objectif en \AnneeEtfBasis, \`a \AlterEtfBasis~ans, avec la m\^eme \'epargne que le
sc\'enario de base, soit \EcartAgeEtfMaisonBasis~ans plus t\^ot. Il part avant l'\^age l\'egal d\`es
\QuoteAnticipeeEtfReferenz~\% d'\'epargne, \EpargneMoisAnticipeeEtfReferenz~€ par mois au d\'ebut, en
\AnneeAnticipeeEtfReferenz, et combin\'e \`a \SparquoteFuenfzig~\% d'\'epargne, il atteint l'objectif \`a
\AlterEtfReferenzFuenfzig~ans. Rejou\'e sur l'histoire am\'ericaine, il a aussi tenu plus souvent que le
revenu de dividendes (section~\ref{sec:retraite-quatre}). Son risque est d'une autre nature~: il vend
des parts chaque ann\'ee, y compris apr\`es un krach, et il a \'epuis\'e le capital des retraites
commenc\'ees en \RetraitHistEchecsListeVierzig.

\vd{Si ton but reste de partir avant l'\^age l\'egal, le calcul de ce document d\'esigne l'ETF monde avec
un retrait de \TauxRetrait~\% comme le levier le plus efficace~; la strat\'egie des dividendes seuls ne
le permet qu'avec au moins \QuoteAnticipeeMaisonSiebzig~\% d'\'epargne.} Ce n'est pas un conseil en
investissement r\'eglement\'e, mais la lecture honn\^ete des chapitres~\ref{sec:capital}
\`a~\ref{sec:retraite}. Choisis ta strat\'egie avant ton premier versement~: en changer plus tard
obligerait \`a vendre des parts et \`a payer l'imp\^ot sur les plus-values r\'ealis\'ees, ce que ce plan
ne chiffre pas.

\subsection{Les premiers pas, d'ao\^ut \`a d\'ecembre 2027}\label{sec:plan-premiers}

\annahme{D\'emarches d\'ecrites avec les r\`egles fiscales de 2026 du chapitre~\ref{sec:fiscalite}~:
leurs modalit\'es pratiques (ordre d'exon\'eration, pr\'el\`evement de la Vorabpauschale, dur\'ee de
validit\'e du formulaire W-8BEN) sont \`a faire confirmer par ton courtier et par un conseiller
fiscal. R\'eserve d'urgence fix\'ee \`a \ReserveUrgenceMois~salaires nets du d\'ebut de carri\`ere,
\GehaltNetEntreeMois~€ chacun, soit \ReserveUrgence~€~: c'est une convention de prudence de ce
document, pas une r\`egle sourc\'ee. Elle est distincte de la r\'eserve de retraite du
chapitre~\ref{sec:retraite}.}

\begin{itemize}
  \item \textbf{Ao\^ut 2027, avant ton premier salaire~: choisir ton courtier.} Cherche un compte-titres
  allemand qui retient lui-m\^eme l'imp\^ot (Abgeltungsteuer, contribution de solidarit\'e,
  Vorabpauschale), qui propose des plans d'\'epargne sur ETF et sur actions et qui g\`ere le
  formulaire W-8BEN. Compare surtout le co\^ut par ordre~: la part maison de ton \'epargne de base,
  \EpargneMoisMaisonBasis~€ par mois, ne repr\'esente que \EpargneMoisParTitre~€ par titre si tu
  ach\`etes chaque mois les \NombreTitres{} titres, et le mod\`ele ne compte aucun frais. Ce document
  ne compare aucun courtier nomm\'ement, faute de source.
  \item \textbf{\`A l'ouverture du compte~: formulaire W-8BEN.} C'est le second \'ecart assum\'e au
  plan~: tous les calculs supposent ce formulaire d\'epos\'e. \PartUs~\% des titres du
  portefeuille-exemple sont am\'ericains, et sans formulaire, cent euros de dividende am\'ericain ne te
  laissent que \NettoUsSansFormulaire~€ au lieu de \NettoUs~€ (section~\ref{sec:fisc-quellensteuer}).
  V\'erifie qu'il est enregistr\'e et demande quand le renouveler.
  \item \textbf{Le m\^eme jour~: ordre d'exon\'eration (Freistellungsauftrag).} Il permet au courtier
  d'appliquer ton forfait d'\'epargnant de \Sparerpauschbetrag~€ par an\QH{sparerpauschbetrag} sans
  attendre ta d\'eclaration~; ce forfait couvre une bonne partie de tes dividendes des premi\`eres
  ann\'ees.
  \item \textbf{Septembre 2027, premier salaire~: r\'eserve d'urgence.} Place les \ReserveUrgence~€ sur
  un compte d'\'epargne disponible, s\'epar\'e du compte-titres. Ton capital de d\'epart y suffit en
  partie~; au milieu de ta fourchette, il manque \ReserveUrgenceComplement~€, soit
  \ReserveUrgenceMoisComplement~mois de ton \'epargne de base.
  \item \textbf{Septembre 2027~: premier versement.} Mets en place un virement de
  \EpargneMoisBasis~€ par mois, \ReferenzSparquote~\% de ton salaire net, ou le montant du taux choisi
  \`a la section~\ref{sec:plan-verdict}, r\'eparti selon ta strat\'egie. Pour limiter les frais, ach\`ete
  les titres de l'ETF maison par roulement, quelques-uns chaque mois, en gardant des poids \'egaux.
\end{itemize}

La r\'eserve d'urgence prend la place de ton capital de d\'epart dans le compte-titres, et c'est un
choix d\'elib\'er\'e. Le chapitre~\ref{sec:accumulation} a montr\'e qu'entre \KapitalDepartNiedrig~€ et
\KapitalDepartHoch~€ de d\'epart, l'\^age d'atteinte ne change pas~: quelques milliers d'euros plac\'es
en 2027 p\`esent peu face \`a des d\'ecennies d'\'epargne, alors qu'une d\'epense impr\'evue pay\'ee en
vendant des parts au mauvais moment co\^uterait davantage. Le calcul n'a toutefois pas \'et\'e refait
pour un capital de d\'epart nul~; l'effet est probablement faible, il n'est pas mesur\'e.

\subsection{Les jalons de ton capital}\label{sec:plan-jalons}

La figure~\ref{fig:jalons-capital-base} place sur une frise les ann\'ees o\`u le capital du sc\'enario
de base franchit quelques seuils, jusqu'\`a l'objectif.

\linienfigur{fig_zeitplan}[width=0.9\linewidth, xlabel={Ann\'ee}, ylabel={Capital (euros de 2026, \'echelle log.)},
  xticklabel style={/pgf/number format/set thousands separator={}}, ymode=log,
  log ticks with fixed point, enlarge y limits={value=0.3}, enlarge x limits=0.08, clip=false,
  table/meta=libelle, point meta=explicit symbolic, nodes near coords,
  every node near coord/.append style={font=\scriptsize, anchor=south, yshift=2pt},
  every axis plot post/.append style={only marks, mark=*, mark size=2pt},
  legend pos=north west, legend cell align=left, legend style={fill=white}]%
  {jahr}{kapital}%
  {{Capital au franchissement du seuil}}%
  {Ann\'ees o\`u le capital du sc\'enario de base franchit \JalonSeuilDixMille, \JalonSeuilCinquanteMille,
  \JalonSeuilCentMille, \JalonSeuilDeuxCentCinquanteMille{} et \JalonSeuilCinqCentMille~€, puis atteint
  l'objectif (r\'epartition \ReferenzStrategie, \ReferenzSparquote~\% d'\'epargne, capital de d\'epart de
  \KapitalDepartMitte~€, euros de 2026, \'echelle logarithmique).\label{fig:jalons-capital-base}}

Ce graphique montre un capital qui passe \JalonSeuilCentMille~€ en \JalonAnneeCentMille{} et
\JalonSeuilCinqCentMille~€ en \JalonAnneeCinqCentMille, puis atteint \KapitalYannBasis~€ en
\AnneeYannBasis. Les premiers seuils tombent vite, \JalonSeuilDixMille~€ d\`es \JalonAnneeDixMille{}
(le mod\`ele compte cette ann\'ee enti\`ere), \JalonSeuilCinquanteMille~€ en
\JalonAnneeCinquanteMille{} et \JalonSeuilDeuxCentCinquanteMille~€ en
\JalonAnneeDeuxCentCinquanteMille~; mais l'essentiel du capital se forme dans les derni\`eres
d\'ecennies, quand les gains d\'epassent tes versements (section~\ref{sec:accu-interets}). Un retard
pris t\^ot se paie donc cher \`a la fin~: chaque ann\'ee d'\'epargne perdue repousse l'objectif d'une
ann\'ee enti\`ere (section~\ref{sec:accu-retard}).

Ces jalons sont en euros de 2026, alors que tes relev\'es de compte seront en euros courants. Pour
comparer ton capital \`a la frise, divise-le par la hausse des prix cumul\'ee depuis 2026~; un capital
en retard de plus d'un an sur la frise annonce un objectif plus tardif.

\subsection{Les points de r\'evision}\label{sec:plan-revision}

Un plan sur plus de quarante ans ne tient que s'il est revu. Ce document se recalcule~: il suffit de
mettre \`a jour les hypoth\`eses sourc\'ees (taux, seuils, salaire) et de relancer le calcul pour
obtenir les m\^emes chapitres avec les valeurs du jour.

\begin{itemize}
  \item \textbf{Chaque janvier.} V\'erifie l'ordre d'exon\'eration et le pr\'el\`evement de la
  Vorabpauschale, compare ton capital, ramen\'e en euros de 2026, \`a la frise de la
  figure~\ref{fig:jalons-capital-base}, et reporte les nouvelles valeurs officielles dans le calcul.
  \item \textbf{\`A chaque augmentation de salaire.} Rel\`eve ton virement pour garder ton taux
  d'\'epargne~: le plan suppose une \'epargne proportionnelle \`a ton salaire net
  (section~\ref{sec:accu-taux}).
  \item \textbf{\`A chaque r\'eforme fiscale ou sociale.} Imp\^ot sur les revenus du capital,
  exon\'eration partielle des ETF ou plafond de l'assurance maladie~: recalcule, car un plafond
  relev\'e repousserait l'objectif (chapitre~\ref{sec:accumulation}).
  \item \textbf{Chaque ann\'ee, pour l'ETF maison.} Rejoue la s\'election du
  chapitre~\ref{sec:portefeuille} avec des donn\'ees fra\^iches~; la liste de ce document n'en est
  qu'une photographie.
  \item \textbf{En \AnneeRevueAvantObjectif, \RevueAvantObjectifAns~ans avant l'objectif.} Revue
  g\'en\'erale~: relev\'e de ta rente l\'egale, assurance maladie des retrait\'es (non mod\'elis\'ee),
  constitution d'une r\'eserve d'un an de d\'epenses, \ZielJahr~€, et choix de tes r\`egles de retrait
  (section~\ref{sec:retraite-regles}). Sur la route de l'ETF monde, cette revue vient
  \RevueAvantObjectifAns~ans avant \AnneeEtfBasis.
\end{itemize}

\bk{Retiens de ce chapitre qu'au taux d'\'epargne de base, ce plan t'am\`ene \`a ton objectif \`a
\AlterYannBasis~ans, l'\^age l\'egal, et que deux d\'ecisions de l'automne 2027 peuvent l'avancer~:
\'epargner au moins \QuoteAnticipeeMaisonSiebzig~\% de ton salaire, ou choisir l'ETF monde avec un
retrait de \TauxRetrait~\%, qui t'y m\`ene \`a \AlterEtfBasis~ans avec la m\^eme \'epargne.} Le reste du
plan tient en quelques gestes~: un courtier, un formulaire W-8BEN, un ordre d'exon\'eration, une
r\'eserve d'urgence, un virement qui suit ton salaire, et une revue chaque ann\'ee.

\Yannbox{\AnneeYannBasis}{En septembre 2027, Yann d\'epose son formulaire W-8BEN et son ordre
d'exon\'eration de \Sparerpauschbetrag~€, met \ReserveUrgence~€ de c\^ot\'e et commence \`a \'epargner
\EpargneMoisBasis~€ par mois. Son capital passe \JalonSeuilCentMille~€ en \JalonAnneeCentMille{} et
\JalonSeuilCinqCentMille~€ en \JalonAnneeCinqCentMille. En \AnneeRevueAvantObjectif, il fait sa revue
g\'en\'erale, et en \AnneeYannBasis, \`a \AlterYannBasis~ans, l'\^age l\'egal, il r\'eunit
\KapitalYannBasis~€. S'il avait \'epargn\'e \QuoteAnticipeeMaisonSiebzig~\% de son salaire, il se
serait arr\^et\'e en \AnneeAnticipeeMaisonSiebzig, \`a \AlterAnticipeeMaisonSiebzig~ans~; avec un ETF
monde et un retrait de \TauxRetrait~\%, d\`es \AnneeEtfBasis, \`a \AlterEtfBasis~ans.}
